\documentclass[12pt]{article}
\usepackage[utf8]{inputenc}
\usepackage[T1]{fontenc}
\usepackage{hyperref}
\title{Southern Lalo Language Resources for K-12 Curriculum Development}
\author{Compiled by Claude Sonnet 5.5 (claude-sonnet-5-5) for Andr\'as Kornai}
\date{}
\begin{document}
\maketitle

\section{Introduction}
Southern Lalo (ISO 639-3 svl; Glottocode sout3210) is a Loloish (Sino-Tibetan) variety of Yunnan, China. In my searches the name appeared only as an alternate name for Mangdi, a Wikipedia entry that gives about 3,000 speakers in Gengma and Cangyuan counties (a 2010 figure). Those pages say an ISO proposal for Mangdi was rejected. I found this only on Wikipedia and its mirrors. I did not confirm it on Glottolog or Ethnologue, so the naming should be checked. I found no school materials, primers, textbooks or native-language resources specific to this variety.

\section{Primary Educational Resources}
I found no primers, textbooks or school programmes for Southern Lalo or Mangdi.

\section{Linguistic and Reference Materials}
\begin{itemize}
\item Mangdi language (Wikipedia), the entry that lists Southern Lalo as an alternate name: \url{https://en.wikipedia.org/wiki/Mangdi_language}
\item Lalo language (Wikipedia), an overview of the wider Lalo cluster and its dialects: \url{https://en.wikipedia.org/wiki/Lalo_language}
\item Susanna Bj\"orverud, \emph{A Grammar of Lalo} (Lund University, 1998). It describes the Longjie (West Mountain) variety, not Southern Lalo, so it is useful only as related-language background: \url{https://www.lunduniversity.lu.se/lup/publication/00a3c596-d8d1-4438-a7fa-abe79ceead0a}
\end{itemize}

\section{Religious and Community Resources}
I found nothing in my searches.

\section{Digital and Audio Resources}
I found nothing in my searches.

\section{Curriculum Development Context}
There is almost no documentation specific to this variety. Any curriculum would need new fieldwork and an orthography decision. Related Lalo descriptions can serve as comparative background only. The Chinese-medium school system is the likely context, which I did not research.

\end{document}